\documentclass[letterpaper,10pt]{article}

% ---------- Packages ----------
\usepackage[T1]{fontenc}
\usepackage{charter}
\usepackage[left=0.3in,right=0.3in,top=0.3in,bottom=0.3in]{geometry}
\usepackage{titlesec}
\usepackage{enumitem}
\usepackage[hidelinks]{hyperref}
\usepackage{xcolor}
\usepackage{tabularx}
\ifdefined\pdfgentounicode
  \input{glyphtounicode}
  \pdfgentounicode=1
\fi

\hypersetup{
  pdftitle={Mujahid Ali Quidwai - Public CV},
  pdfauthor={Mujahid Ali Quidwai}
}

\pagestyle{empty}
\raggedright
\setlength{\parindent}{0pt}
\setlength{\tabcolsep}{0pt}

% ---------- Section style ----------
\titleformat{\section}{\scshape\large\bfseries}{}{0em}{}[\vspace{-2pt}\titlerule]
\titlespacing*{\section}{0pt}{5pt}{2pt}

% ---------- Commands ----------
\newcommand{\entry}[4]{%
  \begin{tabularx}{\textwidth}{@{}X r@{}}
    \textbf{#1} & #2 \\
    \textit{#3} & \textit{#4} \\
  \end{tabularx}\vspace{-3pt}}
\newcommand{\role}[2]{%
  \begin{tabularx}{\textwidth}{@{}X r@{}}
    \textit{#1} & \textit{#2} \\
  \end{tabularx}\vspace{-3pt}}
\newenvironment{bullets}{\begin{itemize}[leftmargin=12pt,itemsep=1pt,parsep=0pt,topsep=1.5pt]}{\end{itemize}\vspace{3pt}}
\newenvironment{subbullets}{\begin{itemize}[leftmargin=12pt,itemsep=0.5pt,parsep=0pt,topsep=1pt,label=$\circ$]}{\end{itemize}}
\newcommand{\lead}[1]{\textbf{#1}}
\newcommand{\centry}[4]{%
  \begin{tabularx}{\textwidth}{@{}X r@{}}
    \textbf{#1}, \textit{#2} & #3 $\cdot$ \textit{#4} \\
  \end{tabularx}\vspace{-3pt}}
\newcommand{\tbd}[1]{\textcolor{red}{\textbf{[#1]}}}


\begin{document}
\begin{center}
  {\LARGE \textbf{Mujahid Ali Quidwai}}\\[3pt]
  \small AI Systems Engineer $\mid$ Agents, Model Post-Training \& Evaluation\\[3pt]
  \href{mailto:quidwaiali@gmail.com}{quidwaiali@gmail.com} $\mid$
  \href{https://www.linkedin.com/in/mujahid-ali-q/}{LinkedIn} $\mid$
  \href{https://scholar.google.com/citations?user=rAAtZi0AAAAJ}{Google Scholar} $\mid$
  \href{https://ali-maq.github.io/}{Portfolio} $\mid$
  \href{https://github.com/Ali-Maq}{GitHub}
\end{center}
\vspace{-7pt}

\section{Profile}
I build AI systems that turn complex data into working tools. My experience spans agent orchestration, model adaptation,
evaluation, and research software, with an emphasis on traceable outputs and comparisons against meaningful baselines.
Healthcare has been a major application area; I also build and evaluate software-agent workflows in public.

\section{Selected Experience}
\entry{Icahn School of Medicine at Mount Sinai}{New York, NY}{Computational Scientist, Oncological Sciences}{Dec 2024 -- Jun 2026}
\begin{bullets}
  \item \lead{OncoCITE:} Designed and implemented a multimodal, six-agent system with 22 MCP tools for extracting structured
  oncology evidence from scientific papers. Linked records to source quotes and page references; used ontology resolution,
  a critic loop, and checkpoints for inspection and resuming longer runs.
  \item \lead{Evidence evaluation:} In 10 retrospective myeloma papers, recovered \textbf{21/25 source-verified reference items};
  expert review confirmed \textbf{44/45 novel extractions} and identified discrepancies in 8/33 curated items. Five additional
  prospective papers were evaluated separately. Published the \href{https://doi.org/10.21203/rs.3.rs-9160944/v1}{preprint}
  and \href{https://github.com/Ali-Maq/civic-extraction-agent}{implementation}; under review at Nature Communications.
  \item \lead{OncoDIF model adaptation:} Fine-tuned a Llama-3 model with LoRA on \textbf{21,000 synthetic question--SQL pairs}.
  Correct SQL improved from \textbf{55\% to 83.24\%} on the study's unseen questions. Work reported in an ASH 2025 conference
  abstract; evaluated terminology variation and flagged low-confidence outputs for review.
  \item \lead{GPU genomics:} Designed a 57-process Nextflow DSL2 pipeline for RNA-seq, whole-exome sequencing, and multi-omics
  integration using NVIDIA Parabricks and RAPIDS. Reduced the observed per-patient workflow from \textbf{7 days to 3 hours (56$\times$)}.
  \item \lead{Collaborative research:} Contributed to myeloma studies on survival modeling, tumor/immune multi-omics integration,
  and single-cell analysis, connecting computational workflows with scientific questions.
\end{bullets}

\role{Associate Computational Scientist, Oncological Sciences}{Oct 2023 -- Dec 2024}
\begin{bullets}
  \item \lead{Clinical RAG:} Built a retrieval-augmented chatbot for myeloma research with source-linked responses and abstention
  when retrieval lacked relevant evidence. Evaluated on an oncologist-curated benchmark against GPT-3.5/GPT-4 baselines
  (medRxiv preprint and IMS 2024 conference abstract).
  \item \lead{Wearables and ML:} Contributed data analysis and predictive modeling to a CAR-T pilot study with 25 evaluable patients,
  published in JCI Insight in 2026. Wearable monitoring detected 18/20 CRS episodes a median of about seven hours before
  nursing detection; this was a separate analysis from the predictive models.
  \item \lead{Domain adaptation:} Fine-tuned Whisper Small for myeloma terminology recognition, with on-device use as the target.
  \item \lead{Mentorship:} Mentored six Carnegie Mellon master's students on CAR-T monitoring research (Sep--Dec 2024), guiding
  experiment design, baseline comparisons, and reporting in collaboration with the research team.
\end{bullets}

\section{Selected Public Project}
\centry{AYA}{Agent workflow evaluation}{\href{https://github.com/Ali-Maq/aya-demo}{github.com/Ali-Maq/aya-demo}}{2026}
\begin{bullets}
  \item Developed the controller and evaluation layer on an Onyx fork, with independent acceptance checks, fixed SQL tasks,
  and recorded tool calls. Published baseline comparisons, verification code, and negative results, including cross-episode
  file access and possible benchmark contamination; credited upstream software and agent-generated components.
\end{bullets}

\section{Education}
\centry{New York University, Tandon}{MS, Computer Engineering (Honors)}{New York, NY}{2021 -- 2023}
\vspace{5pt}
\centry{Dr. A.P.J. Abdul Kalam Technical University}{BTech, Computer Science}{India}{2012 -- 2016}

\newpage
\section{Earlier Experience}
\centry{AYA, NYU Entrepreneurial Institute}{Entrepreneurial Fellow}{New York, NY}{Jul -- Oct 2023}
\begin{bullets}
  \item Developed an education product that turned academic video content into learning questions, connecting speech processing,
  NLP, and product development. Selected for NYU Summer Sprint among 10 of 150 startups. Separate from the 2026 AYA agent project.
\end{bullets}
\centry{New York University}{Machine Learning Researcher \& Teaching Assistant}{New York, NY}{2021 -- 2023}
\begin{bullets}
  \item First-author BEA 2023 workshop paper with Chunhui Li and Parijat Dube on sentence- and document-level AI-text detection;
  achieved \textbf{94\% document-level accuracy on HC3} using paraphrase-generated references and contrastive similarity.
  \item Supported Data Structures \& Algorithms for Bioinformatics under Prof. Manpreet S. Katari for three semesters,
  teaching and assisting more than 130 master's students through assignments, office hours, and project feedback.
\end{bullets}
\centry{Carcrew}{Technical Lead, Founding Team}{New Delhi, India}{Aug 2016 -- Sep 2017}
\begin{bullets}
  \item Built the automotive marketplace MVP with Django, Flask, and PostgreSQL; contributed to technical due diligence for
  a \$2M Series A from TVS Group. The platform grew to more than 10,000 daily active users.
\end{bullets}

\section{Publications}
\begin{bullets}
  \item \href{https://doi.org/10.21203/rs.3.rs-9160944/v1}{\lead{OncoCITE: Multimodal Multi-Agent Reconstruction of Clinical Oncology Knowledge Bases from Scientific Literature.}}
  First author. Research Square preprint, 2026. Under review at Nature Communications.
  \item \href{https://doi.org/10.1172/jci.insight.203988}{\lead{Detection of cytokine release syndrome using wearables and cytokine profiling following CAR-T therapy for myeloma.}}
  Coauthor; data analysis. \textit{JCI Insight} 11(12):e203988, 2026.
  \item \href{https://doi.org/10.1016/s2152-2650(26)03675-x}{\lead{Multi-Modal Single-Cell Profiling of MRD-Negative Patients Identifies Early Immune Signatures Predicting Prolonged Response or Relapse in Multiple Myeloma.}}
  Coauthor. \textit{Clinical Lymphoma Myeloma and Leukemia} 26:S223, 2026. Conference abstract.
  \item \href{https://doi.org/10.1182/blood-2025-2646}{\lead{OncoDIF: An auditable AI framework for automated genomic curation and natural-language clinical querying in hematologic malignancies.}}
  First author. \textit{Blood} 146(Suppl 1):2646, 2025. ASH conference abstract.
  \item \href{https://doi.org/10.1182/blood-2025-3996}{\lead{Development of predictive relapse indicators for myeloma T cell engagers (PRIME) model.}}
  Coauthor. \textit{Blood} 146(Suppl 1):3996, 2025. ASH conference abstract.
  \item \href{https://doi.org/10.1016/s2152-2650(25)03699-7}{\lead{Integrating Microenvironment with Tumor Multi-Omic using Unsupervised Machine Learning to Model Heterogeneity Refines Multiple Myeloma Subtypes and Reveals Immune-Based Clusters with Prognostic Impact.}}
  Coauthor. \textit{Clinical Lymphoma Myeloma and Leukemia} 25:S181--S182, 2025. IMS conference abstract.
  \item \href{https://doi.org/10.1101/2024.03.14.24304293}{\lead{A RAG Chatbot for Precision Medicine of Multiple Myeloma.}}
  First author. medRxiv preprint, 2024.
  \item \href{https://doi.org/10.1016/s2152-2650(24)02048-2}{\lead{Innovative AI-Driven Decision Support Tool for Multiple Myeloma Using Retrieval Augmented Generation.}}
  First author. \textit{Clinical Lymphoma Myeloma and Leukemia} 24:S123--S124, 2024. IMS conference abstract.
  \item \href{https://aclanthology.org/2023.bea-1.58/}{\lead{Beyond Black Box AI generated Plagiarism Detection: From Sentence to Document Level.}}
  First author (published as Ali Quidwai), with Chunhui Li and Parijat Dube. BEA 2023 workshop, ACL Anthology, pp.727--735.
\end{bullets}

\section{Research and Engineering Focus}
Agent orchestration and MCP tools; source-grounded extraction; model adaptation with LoRA; synthetic training data;
baseline comparisons and failure analysis; retrieval-augmented generation; Python, SQL, and reproducible data workflows.

\end{document}
